\documentclass{article}
\usepackage{apacite}
\usepackage{hyperref}
\usepackage{multicol}
\usepackage{natbib}
\usepackage{sophia}
\usepackage[utf8]{inputenc}
\usepackage{array}
\setlength{\parindent}{0pt}
\usepackage{hyperref}
\usepackage{physics}
\usepackage{soul}
\makeatletter
\newcommand*{\rom}[1]{\expandafter\@slowromancap\romannumeral #1@}

\usetikzlibrary{calc,patterns,angles,quotes}
\usepackage{pgfplots}
\pgfplotsset{width=15cm,height=10cm,compat=1.9}

\pgfplotstableread{
x	y	delta-x	delta-y
-0.009	-6.3255	0.0284	0.0028
0.0611	-5.9201	0.0302	0.0019
0.2967	-5.6324	0.0375	0.0014
0.4689	-5.4092	0.0445	0.0011
0.4903	-5.2269	0.0455	0.0009
0.4822	-5.0728	0.0451	0.0008
0.5518	-4.9392	0.0485	0.0007
}{\mytable}

\begin{filecontents}{graphdata.dat}
x	y	delta-x	delta-y
-0.009	-6.3255	0.0284	0.0028
0.0611	-5.9201	0.0302	0.0019
0.2967	-5.6324	0.0375	0.0014
0.4689	-5.4092	0.0445	0.0011
0.4903	-5.2269	0.0455	0.0009
0.4822	-5.0728	0.0451	0.0008
0.5518	-4.9392	0.0485	0.0007
\end{filecontents}

\usepackage[english]{babel}
\usepackage{fancyhdr}
\usepackage{lastpage}
\pagestyle{fancy}
\fancyhf{}


\pagestyle{empty}% for cropping
\usepackage{tikz}
\usetikzlibrary{calc}
\usepackage{zref-savepos}
\usepackage{tabu}

\newcounter{NoTableEntry}
\renewcommand*{\theNoTableEntry}{NTE-\the\value{NoTableEntry}}

\newcommand*{\strike}[2]{%
  \multicolumn{1}{#1}{%
    \stepcounter{NoTableEntry}%
    \vadjust pre{\zsavepos{\theNoTableEntry t}}% top
    \vadjust{\zsavepos{\theNoTableEntry b}}% bottom
    \zsavepos{\theNoTableEntry l}% left
    \hspace{0pt plus 1filll}%
    #2% content
    \hspace{0pt plus 1filll}%
    \zsavepos{\theNoTableEntry r}% right
    \tikz[overlay]{%
      \draw
        let
          \n{llx}={\zposx{\theNoTableEntry l}sp-\zposx{\theNoTableEntry r}sp-\tabcolsep},
          \n{urx}={\tabcolsep},
          \n{lly}={\zposy{\theNoTableEntry b}sp-\zposy{\theNoTableEntry r}sp},
          \n{ury}={\zposy{\theNoTableEntry t}sp-\zposy{\theNoTableEntry r}sp}
        in
        (\n{llx}, \n{lly}) -- (\n{urx}, \n{ury})
      ;
    }% 
  }%
}

\usepackage{tikz}
\usetikzlibrary{calc}

\tikzset{
    right angle quadrant/.code={
        \pgfmathsetmacro\quadranta{{1,1,-1,-1}[#1-1]}     % Arrays for selecting quadrant
        \pgfmathsetmacro\quadrantb{{1,-1,-1,1}[#1-1]}},
    right angle quadrant=1, % Make sure it is set, even if not called explicitly
    right angle length/.code={\def\rightanglelength{#1}},   % Length of symbol
    right angle length=2ex, % Make sure it is set...
    right angle symbol/.style n args={3}{
        insert path={
            let \p0 = ($(#1)!(#3)!(#2)$) in     % Intersection
                let \p1 = ($(\p0)!\quadranta*\rightanglelength!(#3)$), % Point on base line
                \p2 = ($(\p0)!\quadrantb*\rightanglelength!(#2)$) in % Point on perpendicular line
                let \p3 = ($(\p1)+(\p2)-(\p0)$) in  % Corner point of symbol
            (\p1) -- (\p3) -- (\p2)
        }
    }
}
\pagestyle{fancy}
\fancyhf{}
\title{Trigonometric Identities}
\author{Andi Tse}
\date{\today}
\cfoot{Page \thepage \hspace{1pt} of \pageref{LastPage}}
\pagestyle{fancy}



\begin{document}
 \maketitle
 \[a\sin\theta+b\cos\theta=
 \begin{cases}
 \begin{cases}
 
 \sqrt{a^2+b^2}\sin\left(\theta +\arctan\left(\frac ba\right)\right)\quad (a>0)\\
 -\sqrt{a^2+b^2}\sin\left(\theta +\arctan\left(\frac ba\right)\right)\quad (a<0)
 \end{cases}
 \\\\
 \begin{cases}
 
 \sqrt{a^2+b^2}\cos\left(\theta -\arctan\left(\frac ab\right)\right)\quad (b>0)\\
 -\sqrt{a^2+b^2}\cos\left(\theta -\arctan\left(\frac ab\right)\right)\quad (b<0)
 \end{cases}
 \end{cases}
 \]
\begin{align*}
    \sin 2x&=2\sin x\cos x    & \sin\left(\frac x2\right)&=\pm\sqrt{\frac{1-\cos x}{2}}\\
    \cos 2x&=\cos^2x-\sin^2x=1-2\sin^2x=2\cos^2x-1 & \cos\left(\frac x2\right)&=\pm\sqrt{\frac{1+\cos x}{2}}\\
    \tan 2x&=\frac{2\tan x}{1-\tan^2x}&\tan\left(\frac x2\right)&=\frac{1-\cos x}{\sin x}
\end{align*}

\begin{align*}
    \sin(x\pm y)&=\sin x\cos y\pm\cos x\sin y & \sin x\cos y&=\frac 12\left[\sin(x-y)+\sin(x+y)\right]\\
    \cos(x\pm y)&=\cos x\cos y\mp \sin x\sin y&\cos x\cos y&=\frac 12 \left[\cos(x-y)+\cos(x+y)\right]\\
    \tan(x\pm y)&=\frac{\tan x\pm \tan y}{1\mp\tan x\tan y}&
    \sin x\sin y&=\frac 12\left[\cos(x-y)-\cos(x+y)\right]
\end{align*}

\begin{align*}
    \sin x+\sin y&=2\sin\left(\frac{x+y}{2}\right)\cos\left(\frac{x-y}{2}\right) & \sin 3x&=3\sin x-4\sin^3x\\
    \cos x+\cos y&=2\cos\left(\frac{x+y}{2}\right)\cos\left(\frac{x-y}{2}\right)&\cos 3x&=4\cos^3x-3\cos x\\
    \cos x-\cos y&=-2\sin\left(\frac{x+y}{2}\right)\sin\left(\frac{x-y}{2}\right)&\tan 3x&=\frac{3\tan x-\tan^3x}{1-3\tan^2x}
\end{align*}
\end{document}